\begin{abstract}
A language model that is to act, and not only to converse, must state how likely each of its answers is
to be right, and software must be able to act on that number: answer, abstain, or hand the case off.
Calibration research long treated this confidence as something estimated after the fact, from token
probabilities, a post-hoc map, a prompted self-report or repeated sampling. Mostly since 2024, a second line
of work puts calibration into training itself, as a supervised objective or as a reward. This survey maps
both, with the training stage as its core. We collect \nWorks primary works and place each by
the stage at which calibration enters the decision pipeline: the model's own read-off, a post-hoc map,
prompt elicitation, extra inference, or training, plus the works that act on a confidence or fix the
output contract. The \nCore training-stage works, \nCoreRecent of them from 2024 to 2026, are organised
by what the training signal scores, in five families: supervised calibration objectives, reward-model
calibration, calibration rewards on the stated confidence, abstention and refusal training, and
process- and meaning-level confidence rewards. Every work is also classified by the output contract its
confidence is attached to. A typed decision, a probability for each declared
option, is common at the read-off and post-hoc stages, but only \nCoreTyped of the \nCore training works train
one; nearly all the rest train a stated score or a free-text hedge or refusal. Of the
\nSurveysCompared closest prior surveys, none organises calibration training across domains by what the
training signal scores, and none crosses that typology with the output contract. We close with seven
measurable quantities, each tied to a public testbed, and report the first of them, the calibration of
the decision itself, for Jev, a hosted typed-decision model, and \nDIOpen open decision models on a public
leaderboard.
\end{abstract}
